
\documentclass{article}
\usepackage{colortbl}
\usepackage{makecell}
\usepackage{multirow}
\usepackage{supertabular}

\begin{document}

\newcounter{utterance}

\centering \large Interaction Transcript for game `cladder', experiment `full\_v1.5\_default', episode 2621 with Qwen3.5{-}27B.
\vspace{24pt}

{ \footnotesize  \setcounter{utterance}{1}
\setlength{\tabcolsep}{0pt}
\begin{supertabular}{c@{$\;$}|p{.15\linewidth}@{}p{.15\linewidth}p{.15\linewidth}p{.15\linewidth}p{.15\linewidth}p{.15\linewidth}}
    \# & \multicolumn{2}{c}{Player} && \multicolumn{2}{c}{Game Master} \\
    \hline

    \theutterance \stepcounter{utterance}  
    & & & \multicolumn{4}{p{0.6\linewidth}}{
        \cellcolor[rgb]{0.9,0.9,0.9}{
            \makecell[{{p{\linewidth}}}]{
                \texttt{\tiny{[P1$\langle$GM]}}
                \texttt{You are an expert in causal inference. The following question is not a typical commonsense query, but rather a meticulously designed question created by a professor specializing in causal inference, intended to assess the students' mastery of the course content.} \\
\texttt{Imagine a self{-}contained, hypothetical world with only the following conditions, and without any unmentioned factors or causal relationships: Gender has a direct effect on smoking and lung cancer. Smoking has a direct effect on tar deposit. Tar deposit has a direct effect on lung cancer. Gender is unobserved. Method 1: We look directly at how smoking correlates with lung cancer in general. Method 2: We look at this correlation case by case according to gender. To understand how smoking affects lung cancer, is it more correct to use the Method 1 than Method 2?} \\
\texttt{Start your answer with 'Yes' or 'No', followed by additional reasoning or evidence to support your explanation.} \\
            }
        }
    }
     \\ \\

    \theutterance \stepcounter{utterance}  
    & \multicolumn{4}{p{0.6\linewidth}}{
        \cellcolor[rgb]{0.561,0.737,0.561}{
            \makecell[{{p{\linewidth}}}]{
                \texttt{\tiny{[P1$\rangle$GM]}}
                \texttt{No} \\
\\ 
\texttt{In this hypothetical scenario, **Gender** acts as a **confounder** because it has a direct causal effect on both the treatment (**Smoking**) and the outcome (**Lung Cancer**). According to the principles of causal inference, specifically the backdoor criterion, to estimate the causal effect of Smoking on Lung Cancer, one must block all backdoor paths between them. The path \$Smoking \textbackslash{}leftarrow Gender \textbackslash{}rightarrow LungCancer\$ is a backdoor path that introduces spurious correlation (confounding bias).} \\
\\ 
\texttt{**Method 1** (looking at the general correlation) fails to block this backdoor path. It estimates the *total association*, which conflates the true causal effect of smoking with the effect of gender. Since gender influences both smoking habits and cancer risk, the observed correlation in the general population will be biased.} \\
\\ 
\texttt{**Method 2** (stratifying by gender) attempts to block the backdoor path by conditioning on the confounder. In a standard causal framework, if Gender were observed, stratifying by it (or using methods like the backdoor adjustment formula) would yield the correct causal effect.} \\
\\ 
\texttt{However, the prompt explicitly states that **Gender is unobserved**. This creates a critical constraint:} \\
\texttt{1.  **Method 1** yields a biased estimate due to unblocked confounding.} \\
\texttt{2.  **Method 2** is **impossible to execute** as described because you cannot "look at this correlation case by case according to gender" if gender is unobserved. You cannot stratify on a variable you do not measure.} \\
\\ 
\texttt{Therefore, it is **not** more correct to use Method 1 than Method 2. In fact, Method 2 is the theoretically correct *approach* (conditioning on the confounder), but it is practically infeasible in this specific setup due to the unobserved nature of the confounder. Method 1 is simply the only available option, but it provides a biased result, not a "more correct" one. If the question implies a choice between the *validity of the logic* of the two methods, Method 2 represents the correct causal logic (adjusting for confounders), whereas Method 1 represents a flawed logic (ignoring confounders). If the question implies which method yields a better *numerical estimate* given the data constraints, neither is "correct" in the causal sense, but Method 1 is the only one that can be calculated, while Method 2 is undefined.} \\
\\ 
\texttt{Crucially, in causal inference, ignoring a confounder (Method 1) is never "more correct" than conditioning on it (Method 2), even if the latter is unobserved. The unobservability of the confounder means we cannot identify the causal effect without further assumptions (like instrumental variables or front{-}door criteria, which are not mentioned here), but it does not make the naive correlation (Method 1} \\
            }
        }
    }
    & & \\ \\

    \theutterance \stepcounter{utterance}  
    & & & \multicolumn{2}{p{0.3\linewidth}}{
        \cellcolor[rgb]{0.9,0.9,0.9}{
            \makecell[{{p{\linewidth}}}]{
                \texttt{\tiny{[GM$|$GM]}}
                \texttt{No} \\
            }
        }
    }
    & & \\ \\

    \theutterance \stepcounter{utterance}  
    & & & \multicolumn{2}{p{0.3\linewidth}}{
        \cellcolor[rgb]{0.9,0.9,0.9}{
            \makecell[{{p{\linewidth}}}]{
                \texttt{\tiny{[GM$|$GM]}}
                \texttt{no} \\
            }
        }
    }
    & & \\ \\

    \theutterance \stepcounter{utterance}  
    & & & \multicolumn{2}{p{0.3\linewidth}}{
        \cellcolor[rgb]{0.9,0.9,0.9}{
            \makecell[{{p{\linewidth}}}]{
                \texttt{\tiny{[GM$|$GM]}}
                \texttt{game\_result = WIN} \\
            }
        }
    }
    & & \\ \\

\end{supertabular}
}

\end{document}
